\begin{apendicesenv}

\partapendices

\chapter{Síntese da Lean Inception}

A Lean Inception foi usada como protocolo metodológico para definir o MVP do
Avaliador de Serviços Públicos. A atividade reuniu quatro perspectivas: Felipe,
atendimento público; Lucas, gestão digital; Diassis, viabilidade técnica; e
Bruno, experiência do usuário.

\section*{Agenda}

\begin{itemize}
    \item Kick-off e visão do produto;
    \item dinâmica É / Não É / Faz / Não Faz;
    \item objetivos e hipótese do MVP;
    \item personas de cidadão e gestor;
    \item jornada do cidadão no serviço digital;
    \item brainstorming de features contextuais;
    \item revisão técnica, de negócio e UX;
    \item sequenciador;
    \item Canvas MVP.
\end{itemize}

\section*{Escopo Delimitado}

O produto é uma solução de feedback contextual, anônima e agregada, voltada à
melhoria de serviços públicos digitais. Não é ouvidoria, chatbot, sistema de
espionagem, ferramenta de identificação individual ou implantação governamental
completa.

\section*{MVP}

O MVP corresponde às ondas 1, 2 e 3 do sequenciador:

\begin{enumerate}
    \item base técnica de coleta: sessão anônima, cliques, tempo e formulário
    estático comparativo;
    \item detecção e feedback contextual: regra simples, convite contextual,
    estrelas, atributos, NPS e comentário;
    \item painel básico: taxa de sessões com dificuldade, ranking de alertas e
    taxa de resposta dos fluxos.
\end{enumerate}

\chapter{Fluxo Conceitual do Avaliador de Serviços Públicos}

O fluxo conceitual parte do acesso do cidadão a uma página pública ou a uma
página de serviço. O script cliente inicia uma sessão técnica anônima, registra
cliques, página, serviço, origem e instante, e envia os dados para a API. A API
atualiza métricas da sessão e compara cliques e tempo com as faixas esperadas.

Se os sinais estiverem dentro da faixa esperada, a navegação segue normalmente.
Se houver excesso de cliques ou tempo, a API registra dificuldade silenciosa para
o administrador. Em serviços comuns, a navegação pode continuar sem popup. Em
serviços marcados como prioritários ou de alto tráfego, a API pode retornar uma
ação para sugerir redirecionamento.

Quando o usuário apresenta dificuldade persistente ou tempo acima da média, a API
pode retornar ação para abrir uma página simples de feedback. O feedback contém:

\begin{itemize}
    \item avaliação geral por estrelas;
    \item seleção de características do problema, como clareza, tempo, esforço,
    usabilidade, atendimento e eficácia;
    \item pergunta de NPS;
    \item comentário opcional.
\end{itemize}

Ao final, o administrador visualiza alertas, sessões críticas, notas, NPS,
atributos recorrentes e pontos de melhoria.

\chapter{Endpoints e Indicadores do Protótipo}

O protótipo local usa API HTTP e banco SQLite. Os endpoints principais são:

\begin{itemize}
    \item \texttt{GET /health}: verifica disponibilidade;
    \item \texttt{GET /api/services}: lista serviços e regras;
    \item \texttt{POST /api/events/click}: registra clique;
    \item \texttt{POST /api/events/page-time}: registra tempo de permanência;
    \item \texttt{GET /feedback}: exibe formulário de feedback;
    \item \texttt{POST /api/feedback}: salva avaliação estruturada;
    \item \texttt{GET /api/admin/alerts}: lista alertas silenciosos;
    \item \texttt{GET /api/admin/dashboard}: retorna dados agregados.
\end{itemize}

Os indicadores previstos para o painel são:

\begin{itemize}
    \item total de alertas;
    \item sessões com dificuldade;
    \item serviços com maior recorrência de alertas;
    \item média de estrelas por serviço;
    \item NPS por serviço;
    \item taxa de resposta dos fluxos;
    \item atributos mais selecionados;
    \item comentários recentes.
\end{itemize}

\end{apendicesenv}
